% 参考文献。注意：main.tex 中 \referencescn 已通过 thebibliography 自动生成
% “参考文献”一级标题（不带编号），并已 \addcontentsline 进入目录，
% 所以本文件**不要再写一级标题**，否则会出现两个“参考文献”标题。

\begin{thebibliography}{99}
\bibitem{ref1} Lo Y M D, Corbetta N, Chamberlain P F, et al. Presence of fetal DNA in maternal plasma and serum[J]. The Lancet, 1997, 350(9076): 485--487.
\bibitem{ref2} Ashoor G, Syngelaki A, Poon L C Y, et al. Fetal fraction in maternal plasma cell-free DNA at 11--13 weeks' gestation: relation to maternal and fetal characteristics[J]. Ultrasound in Obstetrics \& Gynecology, 2013, 41(1): 26--32.
\bibitem{ref3} Norton M E, Jacobsson B, Swamy G K, et al. Cell-free DNA analysis for noninvasive examination of trisomy[J]. New England Journal of Medicine, 2015, 372(17): 1589--1597.
\bibitem{ref4} Wang E, Batey A, Struble C, et al. Gestational age and maternal weight effects on fetal cell-free DNA in maternal plasma[J]. Prenatal Diagnosis, 2013, 33(7): 662--666.
\bibitem{ref5} Bianchi D W, Chiu R W K. Sequencing of circulating cell-free DNA during pregnancy[J]. New England Journal of Medicine, 2018, 379(5): 464--473.
\bibitem{ref6} 李金明. 无创产前检测技术的临床应用与质量控制[J]. 中华检验医学杂志, 2016, 39(5): 321--324.
\end{thebibliography}
